% GENERATED FILE. Edit canonical lessons, then run npm run generate:books.
% canonical-source-hash: fnv1a64:ccc298c2b1b1720e
% canonical-lessons: HI-C524-galana, HI-C524-parosna, HI-C524-ghotna, HI-C524-jamana, HI-C524-khagalna

\chapter{Cooking and Serving}
\label{ch:hi-cooking-and-serving}
\hlchaptermodality{524}

\begin{chapteropening}
\textbf{By the end of this chapter:} \emph{I can serve a meal and clear up after it.}

Complete the last lesson of chapter 524: I can serve a meal and clear up after it.
\end{chapteropening}

\section[\texorpdfstring{\hi{गलाना} (galānā)}{गलाना (galānā)}]{\hi{गलाना} (galānā) — to cook soft; to soften}

\label{lesson:HI-C524-galana}



\begin{quote}
Before the new one: say the Hindi for to sprinkle (a powder), then the Hindi for to grind coarsely.
\end{quote}

\subsection*{You'll want to know: \hi{गलाना}}
\textbf{\hi{गलाना}} — \emph{galānā} — \textquotedblleft{}to cook soft; to soften\textquotedblright{}.

\begin{quote}
\hi{दाल} \hi{को} \hi{पूरी} \hi{तरह} \hi{गलाइए।}
\end{quote}

\begin{quote}
\emph{Dāl ko pūrī tarah galāie.}
\end{quote}

\begin{quote}
Cook the dal until it is completely soft.
\end{quote}

Past, with \hi{ने} (he or she …): \hi{उसने} \hi{चावल} \hi{गलाए}, \emph{usne cāval galāe}.

\begin{cousinweb}
\textbf{\hi{गलाना}} (\emph{galānā}) is the causative of \emph{galnā}, \textquotedblleft{}to melt, soften\textquotedblright{}, which CDIAL files under Sanskrit \emph{galati}, \textquotedblleft{}drips, oozes, trickles\textquotedblright{} (CDIAL 4074); CDIAL gives Hindi \emph{galnā} as \textquotedblleft{}to melt, trickle\textquotedblright{}.
\end{cousinweb}

\begin{grammarlens}[title={What you've built}]
One more word for describing what you lived and felt.
\end{grammarlens}

\subsection*{Your turn}
\begin{itemize}
\raggedright
  \item \emph{Say it:} \emph{galānā}
  \item \emph{Say it:} \emph{galānā}, once more
  \item \emph{Say it:} \emph{Dāl ko pūrī tarah galāie.} — Cook the dal until it is completely soft.
  \item \emph{Recall:} say the Hindi for to sprinkle (a powder), then the Hindi for to grind coarsely, then say \emph{galānā} again
\end{itemize}

\begin{tcolorbox}[breakable,colback=teal!4,colframe=teal!35!black,title={Before you move on}]
What does \hi{गलाना} mean? (\textquotedblleft{}To cook soft; to soften\textquotedblright{}.) Say it once more.
\end{tcolorbox}
\section[\texorpdfstring{\hi{परोसना} (parosnā)}{परोसना (parosnā)}]{\hi{परोसना} (parosnā) — to serve (food)}

\label{lesson:HI-C524-parosna}



\begin{quote}
Before the new one: say the Hindi for to grind coarsely, then the Hindi for to cook soft.
\end{quote}

\subsection*{You'll want to know: \hi{परोसना}}
\textbf{\hi{परोसना}} — \emph{parosnā} — \textquotedblleft{}to serve (food)\textquotedblright{}.

\begin{quote}
\hi{माँ} \hi{ने} \hi{मेहमानों} \hi{को} \hi{गरम} \hi{खाना} \hi{परोसा।}
\end{quote}

\begin{quote}
\emph{Māṁ ne mehmānoṁ ko garam khānā parosā.}
\end{quote}

\begin{quote}
Mother served the guests hot food.
\end{quote}

Past, with \hi{ने} (he or she …): \hi{उसने} \hi{दाल} \hi{परोसी}, \emph{usne dāl parosī}.

\begin{cousinweb}
\textbf{\hi{परोसना}} (\emph{parosnā}) is a tadbhava, inherited from Sanskrit \emph{parivēṣati}, \textquotedblleft{}serves food\textquotedblright{} (CDIAL 7888).
\end{cousinweb}

\begin{grammarlens}[title={What you've built}]
One more word for describing what you lived and felt.
\end{grammarlens}

\subsection*{Your turn}
\begin{itemize}
\raggedright
  \item \emph{Say it:} \emph{parosnā}
  \item \emph{Say it:} \emph{parosnā}, once more
  \item \emph{Say it:} \emph{Māṁ ne mehmānoṁ ko garam khānā parosā.} — Mother served the guests hot food.
  \item \emph{Recall:} say the Hindi for to grind coarsely, then the Hindi for to cook soft, then say \emph{parosnā} again
\end{itemize}

\begin{tcolorbox}[breakable,colback=teal!4,colframe=teal!35!black,title={Before you move on}]
What does \hi{परोसना} mean? (\textquotedblleft{}To serve (food)\textquotedblright{}.) Say it once more.
\end{tcolorbox}
\section[\texorpdfstring{\hi{घोंटना} (ghõṭnā)}{घोंटना (ghõṭnā)}]{\hi{घोंटना} (ghõṭnā) — to grind into a paste; to rub smooth}

\label{lesson:HI-C524-ghotna}



\begin{quote}
Before the new one: say the Hindi for to cook soft, then the Hindi for to serve (food).
\end{quote}

\subsection*{You'll want to know: \hi{घोंटना}}
\textbf{\hi{घोंटना}} — \emph{ghõṭnā} — \textquotedblleft{}to grind into a paste; to rub smooth\textquotedblright{}.

\begin{quote}
\hi{माँ} \hi{पत्थर} \hi{पर} \hi{चटनी} \hi{घोंट} \hi{रही} \hi{है।}
\end{quote}

\begin{quote}
\emph{Māṁ patthar par caṭnī ghõṭ rahī hai.}
\end{quote}

\begin{quote}
Mother is grinding chutney on a stone.
\end{quote}

Past, with \hi{ने} (he or she …): \hi{उसने} \hi{मसाला} \hi{घोंटा}, \emph{usne masālā ghõṭā}.

\begin{cousinweb}
\textbf{\hi{घोंटना}} (\emph{ghõṭnā}) is a tadbhava, inherited from Sanskrit \emph{ghaṭṭayati}, \textquotedblleft{}rubs\textquotedblright{} (CDIAL 4417). CDIAL gives the Hindi verb as \textquotedblleft{}to rub\textquotedblright{}.
\end{cousinweb}

\begin{grammarlens}[title={What you've built}]
One more word for describing what you lived and felt.
\end{grammarlens}

\subsection*{Your turn}
\begin{itemize}
\raggedright
  \item \emph{Say it:} \emph{ghõṭnā}
  \item \emph{Say it:} \emph{ghõṭnā}, once more
  \item \emph{Say it:} \emph{Māṁ patthar par caṭnī ghõṭ rahī hai.} — Mother is grinding chutney on a stone.
  \item \emph{Recall:} say the Hindi for to cook soft, then the Hindi for to serve (food), then say \emph{ghõṭnā} again
\end{itemize}

\begin{tcolorbox}[breakable,colback=teal!4,colframe=teal!35!black,title={Before you move on}]
What does \hi{घोंटना} mean? (\textquotedblleft{}To grind into a paste; to rub smooth\textquotedblright{}.) Say it once more.
\end{tcolorbox}
\section[\texorpdfstring{\hi{जमाना} (jamānā)}{जमाना (jamānā)}]{\hi{जमाना} (jamānā) — to set (curd); to make (something) set}

\label{lesson:HI-C524-jamana}



\begin{quote}
Before the new one: say the Hindi for to serve (food), then the Hindi for to grind into a paste.
\end{quote}

\subsection*{You'll want to know: \hi{जमाना}}
\textbf{\hi{जमाना}} — \emph{jamānā} — \textquotedblleft{}to set (curd); to make (something) set\textquotedblright{}.

\begin{quote}
\hi{रात} \hi{को} \hi{माँ} \hi{ने} \hi{दूध} \hi{में} \hi{जामन} \hi{डालकर} \hi{दही} \hi{जमाया।}
\end{quote}

\begin{quote}
\emph{Rāt ko māṁ ne dūdh meṁ jāman ḍālkar dahī jamāyā.}
\end{quote}

\begin{quote}
At night my mother added starter to the milk and set the yoghurt.
\end{quote}

Past, with \hi{ने} (he or she …): \hi{उसने} \hi{बर्फ़} \hi{जमाई}, \emph{usne barf jamāī}.

\begin{cousinweb}
\textbf{\hi{जमाना}} (\emph{jamānā}) is a tadbhava, inherited from Sanskrit \emph{yamyatē}, \textquotedblleft{}is held, is fixed\textquotedblright{} (CDIAL 10428). CDIAL lists \emph{jamānā}, \textquotedblleft{}to collect\textquotedblright{}, beside the intransitive \emph{jamnā}.
\end{cousinweb}

\begin{grammarlens}[title={What you've built}]
One more word for describing what you lived and felt.
\end{grammarlens}

\subsection*{Your turn}
\begin{itemize}
\raggedright
  \item \emph{Say it:} \emph{jamānā}
  \item \emph{Say it:} \emph{jamānā}, once more
  \item \emph{Say it:} \emph{Rāt ko māṁ ne dūdh meṁ jāman ḍālkar dahī jamāyā.} — At night my mother added starter to the milk and set the yoghurt.
  \item \emph{Recall:} say the Hindi for to serve (food), then the Hindi for to grind into a paste, then say \emph{jamānā} again
\end{itemize}

\begin{tcolorbox}[breakable,colback=teal!4,colframe=teal!35!black,title={Before you move on}]
What does \hi{जमाना} mean? (\textquotedblleft{}To set (curd); to make (something) set\textquotedblright{}.) Say it once more.
\end{tcolorbox}
\section[\texorpdfstring{\hi{खँगालना} (khãgālnā)}{खँगालना (khãgālnā)}]{\hi{खँगालना} (khãgālnā) — to rinse out}

\label{lesson:HI-C524-khagalna}



\begin{quote}
Before the new one: say the Hindi for to grind into a paste, then the Hindi for to set (curd).
\end{quote}

\subsection*{You'll want to know: \hi{खँगालना}}
\textbf{\hi{खँगालना}} — \emph{khãgālnā} — \textquotedblleft{}to rinse out\textquotedblright{}.

\begin{quote}
\hi{खाने} \hi{के} \hi{बाद} \hi{उसने} \hi{गिलास} \hi{पानी} \hi{से} \hi{खँगाले।}
\end{quote}

\begin{quote}
\emph{Khāne ke bād usne gilās pānī se khãgāle.}
\end{quote}

\begin{quote}
After the meal he rinsed the tumblers with water.
\end{quote}

Past, with \hi{ने} (he or she …): \hi{उसने} \hi{बोतल} \hi{खँगाली}, \emph{usne botal khãgālī}.

\begin{cousinweb}
\textbf{\hi{खँगालना}} (\emph{khãgālnā}) is a tadbhava, inherited from \emph{*khakkhālayati}, \textquotedblleft{}rinses\textquotedblright{}, a form CDIAL reconstructs from its descendants (CDIAL 3762).
\end{cousinweb}

\begin{grammarlens}[title={What you've built}]
One more word for describing what you lived and felt.
\end{grammarlens}

\subsection*{Your turn}
\begin{itemize}
\raggedright
  \item \emph{Say it:} \emph{khãgālnā}
  \item \emph{Say it:} \emph{khãgālnā}, once more
  \item \emph{Say it:} \emph{Khāne ke bād usne gilās pānī se khãgāle.} — After the meal he rinsed the tumblers with water.
  \item \emph{Recall:} say the Hindi for to grind into a paste, then the Hindi for to set (curd), then say \emph{khãgālnā} again
\end{itemize}

\begin{tcolorbox}[breakable,colback=teal!4,colframe=teal!35!black,title={Before you move on}]
What does \hi{खँगालना} mean? (\textquotedblleft{}To rinse out\textquotedblright{}.) Say it once more.
\end{tcolorbox}
